\chapter{Circuits with measurements}
\label{ch:measurement}

A circuit of Chapter~\ref{ch:circuits} may measure sites in the computational basis and choose
its later layers by the outcomes. The measurement circuits of this chapter have unitary gates
and onsite families that are unitaries or measurements; the outcomes are communicated to every
later layer at no cost, the measurements and the onsite unitaries are free, and only the layers
using a cell of positive cost are counted in the depth, as in the model of local circuits with
classical communication of~\cite{Piroli2021LOCC}. A circuit implements a matrix on a set of
vectors when, whatever the outcomes, it acts there as that matrix up to a scalar, and a vector
is converted deterministically into another when every outcome gives a multiple of the second.
Conversions compose with additive depths, and a deterministic conversion of unit vectors is an
adaptive conversion of their density matrices in the sense of Chapter~\ref{ch:protocols}.

\section{Measurement circuits and implementations}
\label{sec:measurement_model}

\begin{definition}[Measurement circuits]
    \label{def:measurement_circuit}
    \lean{KrausFamily.IsUnitaryOrMeasurement,
        KrausFamily.IsUnitary.isUnitaryOrMeasurement,
        KrausFamily.isUnitaryOrMeasurement_measure}
    \leanok
    \uses{def:circuits_circuit, def:circuits_local, def:circuits_outcome}
    A Kraus family from $\C^\alpha$ to $\C^\beta$ is a local family of a measurement circuit if
    it consists of one unitary operator, or if its operators are $\ket{f(a)}\bra a$ for the
    elements $a$ of $\alpha$, indexed by its outcomes through a bijection, for a bijection
    $f\colon\alpha\to\beta$: the measurement of the basis followed by a relabelling. A
    measurement circuit is a circuit of plain gates all of whose gates and onsite families are
    local families of a measurement circuit. Unitary families and the measurement of a basis
    are local families of a measurement circuit.
\end{definition}

\begin{lemma}[Measurement circuits preserve the trace]
    \label{lem:measurement_tp}
    \lean{KrausFamily.IsUnitaryOrMeasurement.isTP}
    \leanok
    \uses{def:measurement_circuit}
    Every local family of a measurement circuit is trace-preserving. Hence every measurement
    circuit is an adaptive protocol of its depth.
\end{lemma}
\begin{proof}
    \leanok
    A unitary $V$ has $V^\dagger V=\mathbf{1}$. For the operators $K_a=\ket{f(a)}\bra a$,
    $K_a^\dagger K_a=\ket a\bra a$ and $\sum_a\ket a\bra a=\mathbf{1}$.
\end{proof}

\begin{definition}[Implementation of a matrix]
    \label{def:measurement_implementation}
    \lean{QuantumCircuit.Circuit.IsImplementationOn}
    \leanok
    \uses{def:circuits_circuit}
    A circuit $c$ from $X$ to $Z$ implements a matrix $W$ from $\mathcal H_X$ to $\mathcal H_Z$
    on a set $E$ of vectors if for every branch $p$ of $c$ and every $v\in E$ there is
    $z\in\C$ with $K_p v=z\,Wv$, where $K_p$ is the Kraus operator of the branch.
\end{definition}

\begin{lemma}[Implementations compose]
    \label{lem:measurement_implementation_comp}
    \lean{QuantumCircuit.Circuit.IsImplementationOn.mono,
        QuantumCircuit.Circuit.isImplementationOn_done,
        QuantumCircuit.Circuit.isImplementationOn_of_kraus_eq,
        QuantumCircuit.Circuit.IsImplementationOn.append,
        QuantumCircuit.Circuit.IsImplementationOn.comp}
    \leanok
    \uses{def:measurement_implementation, thm:circuits_composition}
    An implementation on $E$ is one on every subset of $E$. The empty circuit implements
    $\mathbf{1}$, and a circuit all of whose Kraus operators are $U$ implements $U$ on every set. If
    $c$ implements $W$ on $E$, $W$ maps $E$ into $E'$, and the continuation attached to every
    branch of $c$ implements $W'$ on $E'$, then the composite circuit implements $W'W$ on $E$;
    in particular so does $c$ followed by one circuit implementing $W'$ on $E'$.
\end{lemma}
\begin{proof}
    \leanok
    A branch of the composite is a branch $p$ of $c$ followed by a branch $q$ of its
    continuation, with Kraus operator $K'_qK_p$. For $v\in E$, $K_pv=zWv$ with $Wv\in E'$, so
    $K'_qK_pv=z\,K'_qWv=zz'\,W'Wv$.
\end{proof}

\begin{lemma}[Branch scalars on a subspace]
    \label{lem:measurement_branch_scalar}
    \lean{QuantumCircuit.Circuit.IsImplementationOn.exists_smul_of_isometry}
    \leanok
    \uses{def:measurement_implementation}
    Let $c$ implement an isometry $W$ on a linear subspace $E$. Then for every branch $p$ there
    is one $z\in\C$, possibly zero, with $K_pv=z\,Wv$ for all $v\in E$.
\end{lemma}
\begin{proof}
    \leanok
    The map $v\mapsto W^\dagger K_pv$ sends $E$ to $E$, since $K_pv=z_vWv$ and $W^\dagger W=\mathbf{1}$
    give $W^\dagger K_pv=z_vv$. A linear endomorphism under which every vector is an eigenvector
    is a scalar $z$, and for $v\neq0$ the equation $z_vv=zv$ gives $z_v=z$.
\end{proof}

\begin{lemma}[A nonzero branch]
    \label{lem:measurement_nonzero_branch}
    \lean{QuantumCircuit.Circuit.exists_kraus_mulVec_ne_zero}
    \leanok
    \uses{def:circuits_circuit, thm:circuits_tp}
    If every gate and onsite family of a circuit is trace-preserving, then for every nonzero
    vector $v$ some branch $p$ has $K_pv\neq0$.
\end{lemma}
\begin{proof}
    \leanok
    The circuit is trace-preserving, so $\sum_pK_p^\dagger K_p=\mathbf{1}$; if every $K_pv$ vanished,
    $v=\sum_pK_p^\dagger K_pv=0$.
\end{proof}

\begin{definition}[Measurement with corrections]
    \label{def:measurement_correct}
    \lean{QuantumCircuit.Circuit.measureCorrect,
        QuantumCircuit.Circuit.kraus_measureCorrect,
        QuantumCircuit.Circuit.depth_measureCorrect,
        QuantumCircuit.Circuit.forall_measureCorrect}
    \leanok
    \uses{def:circuits_outcome, lem:circuits_measure_kraus, def:measurement_circuit}
    For a region $D$ and matrices $V_{m,s}$ on $\C^{X_s}$ indexed by the configurations $m$ of
    $D$ and the sites $s$, the measurement of $D$ with the corrections $V$ is the circuit that
    measures $D$ in the computational basis and then applies the onsite layer
    $\bigotimes_sV_{m,s}$ chosen by the outcome $m$. Its Kraus operator on a branch with outcome
    $m$ is $(\bigotimes_sV_{m,s})P_m$, with $P_m$ the projection onto $\ket m$ placed on $D$;
    its depth is $0$, and it is a measurement circuit when the $V_{m,s}$ are unitary.
\end{definition}

\begin{lemma}[A round implementing a matrix]
    \label{lem:measurement_round_implementation}
    \lean{QuantumCircuit.Circuit.isImplementationOn_comp_measureCorrect}
    \leanok
    \uses{def:measurement_correct, def:measurement_implementation}
    Let every Kraus operator of the circuit $u$ be $U$. If for every configuration $m$ of $D$
    there is $z\in\C$ with $(\bigotimes_sV_{m,s})P_mUv=z\,Wv$ for every $v\in E$, then $u$
    followed by the measurement of $D$ with the corrections $V$ implements $W$ on $E$.
\end{lemma}
\begin{proof}
    \leanok
    A branch of the composite has Kraus operator $(\bigotimes_sV_{m,s})P_mU$ for its outcome
    $m$.
\end{proof}

\section{Deterministic conversions}
\label{sec:measurement_conversion}

\begin{definition}[Deterministic conversion with measurements]
    \label{def:measurement_conversion}
    \lean{QuantumCircuit.IsMeasurementConversion}
    \leanok
    \uses{def:measurement_circuit, def:circuits_costs}
    A nonzero vector $\varphi\in\mathcal H_X$ is converted into $\psi\in\mathcal H_Z$ in depth
    $T$ if some measurement circuit from $X$ to $Z$ of depth at most $T$ has, on every branch
    $p$, $K_p\varphi=z_p\psi$ for some $z_p\in\C$. A preparation of $\psi$ is a conversion from
    a basis configuration or from a product vector. The number of rounds of measurements is not
    bounded, whereas~\cite{Piroli2021LOCC} performs one round in its class
    $\mathrm{QCcc}_\ell$ and notes a more general scheme with multiple rounds.
\end{definition}

\begin{theorem}[Conversions compose]
    \label{thm:measurement_conversion_comp}
    \lean{QuantumCircuit.IsMeasurementConversion.ne_zero,
        QuantumCircuit.IsMeasurementConversion.mono,
        QuantumCircuit.IsMeasurementConversion.refl,
        QuantumCircuit.IsMeasurementConversion.smul,
        QuantumCircuit.IsMeasurementConversion.trans,
        QuantumCircuit.Circuit.IsImplementationOn.isMeasurementConversion,
        QuantumCircuit.IsMeasurementConversion.of_prodOp,
        QuantumCircuit.IsUnitaryCircuit.isMeasurementConversion,
        QuantumCircuit.IsPreparedInDepth.exists_isMeasurementConversion}
    \leanok
    \uses{def:measurement_conversion, lem:measurement_implementation_comp,
        lem:measurement_nonzero_branch, lem:measurement_tp}
    A vector converted from a nonzero vector is nonzero. A conversion in depth $T$ is one in
    every depth $T'\ge T$, every nonzero vector is converted into itself in depth $0$, and a
    conversion into $\psi$ is one into $a\psi$ for $a\neq0$. Converting $\varphi$ into $\psi$ in
    depth $T_1$ and $\psi$ into $\chi$ in depth $T_2$ converts $\varphi$ into $\chi$ in depth
    $T_1+T_2$. A measurement circuit of depth at most $T$ implementing $W$ on a set containing
    the nonzero vector $\varphi$ converts $\varphi$ into $W\varphi$ in depth $T$. A conversion
    from $(\bigotimes_sU_s)\varphi$, for unitaries $U_s$, is a conversion from $\varphi$ in the
    same depth; in particular a conversion from a product vector is one from a basis
    configuration. A unitary circuit of depth at most $T$ with matrix $U$ converts every
    nonzero $\varphi$ into $U\varphi$ in depth $T$; hence a nonzero vector prepared in depth $T$
    by a unitary circuit is converted in depth $T$ from a product vector.
\end{theorem}
\begin{proof}
    \leanok
    By Lemma~\ref{lem:measurement_nonzero_branch} some branch has $K_p\varphi\neq0$, and
    $K_p\varphi=z_p\psi$ forces $\psi\neq0$. For the composition, the branch $(p,q)$ of the
    composite circuit gives $K'_qK_p\varphi=z_pK'_q\psi=z_pz'_q\chi$, and depths add by
    Theorem~\ref{thm:circuits_cost_subadditive}. The onsite unitaries form a layer of depth
    $0$ placed before the circuit.
\end{proof}

\begin{theorem}[Deterministic conversions are adaptive conversions]
    \label{thm:measurement_adaptive}
    \lean{QuantumCircuit.IsMeasurementConversion.isAdaptiveConversion}
    \leanok
    \uses{def:measurement_conversion, def:protocols_protocol, lem:measurement_tp}
    If the unit vector $\varphi$ is converted into the unit vector $\psi$ in depth $T$, then
    $\ket\varphi\bra\varphi$ is converted adaptively into $\ket\psi\bra\psi$ in depth $T$.
\end{theorem}
\begin{proof}
    \leanok
    The map of the circuit sends $\ket\varphi\bra\varphi$ to
    $\sum_pK_p\ket\varphi\bra\varphi K_p^\dagger=\bigl(\sum_p|z_p|^2\bigr)\ket\psi\bra\psi$,
    and since the circuit preserves the trace, $\sum_p|z_p|^2=1$.
\end{proof}

\begin{definition}[One round of measurements]
    \label{def:measurement_one_round}
    \lean{QuantumCircuit.IsOneRoundConversion}
    \leanok
    \uses{def:measurement_correct, def:circuits_local}
    A nonzero vector $\varphi$ is converted into $\psi$ with one round of measurements in
    depths $(T_1,T_2)$ if there are a unitary circuit with matrix $U$ of depth at most $T_1$, a
    region $D$, unitaries $V_{m,s}$, and a unitary circuit with matrix $U'$ of depth at most
    $T_2$, such that for every configuration $m$ of $D$,
    $U'(\bigotimes_sV_{m,s})P_mU\varphi$ is a scalar multiple of $\psi$. With $T_2=0$ this is
    the class $\mathrm{QCcc}_{T_1}$ of~\cite{Piroli2021LOCC}, the circuit of depth $0$ being
    free local unitaries: a local circuit, measurements, and local unitaries chosen from all the
    outcomes. With $T_2>0$ it is the composition of such a transformation with a circuit
    without measurements, as in the preparations of~\cite{Malz2024Preparation}. Any site may be
    measured, in the computational basis, and the corrections are onsite unitaries applied
    once after the measurement.
\end{definition}

\begin{lemma}[One round]
    \label{lem:measurement_one_round}
    \lean{QuantumCircuit.IsOneRoundConversion.isMeasurementConversion,
        QuantumCircuit.IsOneRoundConversion.mono,
        QuantumCircuit.IsOneRoundConversion.mulVec,
        QuantumCircuit.IsUnitaryCircuit.isOneRoundConversion}
    \leanok
    \uses{def:measurement_one_round, def:measurement_conversion,
        thm:measurement_conversion_comp}
    A conversion with one round of measurements in depths $(T_1,T_2)$ is a deterministic
    conversion in depth $T_1+T_2$, and one in all larger depths. Following it by a unitary
    circuit of depth at most $T$ with matrix $W$ gives a conversion into $W\psi$ in depths
    $(T_1,T_2+T)$. A unitary circuit of depth at most $T$ with matrix $U$ converts every
    nonzero $\varphi$ into $U\varphi$ with one round in depths $(T,0)$, measuring no site.
\end{lemma}
\begin{proof}
    \leanok
    The first unitary circuit, the measurement of $D$ with the corrections $V$, and the second
    unitary circuit form a measurement circuit whose Kraus operator on a branch with outcome
    $m$ is $U'(\bigotimes_sV_{m,s})P_mU$, of depth at most $T_1+0+T_2$. Measuring the empty
    region gives $P_m=\mathbf{1}$.
\end{proof}
